\documentclass[11pt]{article}
\usepackage[margin=1in]{geometry}
\usepackage{amsmath,amssymb,booktabs,longtable,tabularx,array}
\usepackage{enumitem}
\usepackage[T1]{fontenc}
\usepackage[hidelinks]{hyperref}
\setlist[itemize]{leftmargin=*}
\setlist[enumerate]{leftmargin=*}
\newcommand{\code}[1]{\texttt{\detokenize{#1}}}
\newcommand{\st}{s_t}
\newcommand{\snext}{s_{t+1}}
\newcommand{\obs}{o_t}
\newcommand{\sobs}{\sigma_{\mathrm{obs}}}
\newcommand{\ssafe}{s_{\mathrm{safe}}}
\newcommand{\E}{\mathbb{E}}
\newcommand{\Var}{\mathrm{Var}}
\newcommand{\1}{\mathbf{1}}
\sloppy

\usepackage{graphicx}

\title{Motivation Experiment Results: Native Ecological Baselines}
\author{Ecology Benchmark Documentation}
\date{July 2026}

\begin{document}
\maketitle

\section{Question and Answer}

The motivation experiment asked whether general offline model-based RL methods
beat native discrete ecological baselines in the real-ecology benchmark once
the baselines are implemented as genuine discrete-state planners rather than as
the adapted particle-MPC methods. The expected story was that structural
uncertainty and partial observability would make naive ecological solvers
fragile.

The completed run shows the opposite. The best native ecological solver beats
or ties the general methods in 860 of 864 paired comparisons. \code{refplan}
and \code{bamcts} never beat both native solvers; \code{ogsrl} does so in only
4 of 288 paired cells. More importantly, the result survives without the
stronger \code{plus_native} model-identifying solver: the single-Ricker
\code{moor_native} baseline beats the best general method by 0.819 mean return
on recoverable safe cells, and the general methods win only 6 of 112 such
cells. The result is therefore a negative result for the original motivation,
not a marginally inconclusive result.

\section{Run Scope and Integrity}

The main run contains 1440 completed evaluation summaries:
9 populations, 4 dynamics families, 4 observation-noise values, 2 reward modes,
and 5 methods. The resolution replicate contains 1152 native-only summaries
at coarse and fine native grids. The launch gates passed before the sweep:
the 288-row acceptance battery passed all 8 checks; all 288 main cells share a
single dataset hash across methods; and the native-grid resolution check has
0 headline flips across 288 cells.

\begin{center}
\begin{tabular}{lrrr}
\toprule
artifact & rows & purpose & result \\
\midrule
main sweep & 1440 & headline comparison & complete \\
coarse/fine native grids & 1152 & discretization sensitivity & 0 flips \\
acceptance battery & 288 & launch gate & 8/8 pass \\
search-budget audit & 1728 & under-tuning check & model-limited \\
\bottomrule
\end{tabular}
\end{center}

\section{Benchmark Setting}

The benchmark is a partially observed conservation-control problem. The hidden
state is true population abundance. The policy observes a noisy abundance
survey, public management controls, and the public action history, but it does
not observe the private true state during deployment. The motivation sweep uses
four observation-noise values, $\sobs \in \{0,0.1,0.2,0.4\}$.

Each cell is a combination of one of 9 real populations, one of 4 ecological
dynamics families, one reward mode, and one observation-noise value. The four
dynamics families are Ricker, Allee, theta-logistic, and regime-switching. The
action set has 11 real management actions, including do-nothing, harvest,
restoration, integrated conservation, and translocation-style interventions.
The action table publicly determines action-specific growth set-points and
carrying-capacity effects; the species table publicly determines the
population-specific carrying-capacity scale.

The main comparison uses two reward modes. The yield mode removes the collapse
penalty and behaves like a production baseline. The safe mode uses an
occupancy-style safety penalty with collapse penalty $P=5$, so populations are
penalized while they remain below the cell's safety threshold. Each public
offline dataset contains 4000 logged transitions. Evaluation uses 5 seed blocks,
4 episodes per seed, and horizon 50.

This public structure is the central reason the result is surprising. The
benchmark was intended to stress ecological baselines with partial observability
and structural uncertainty, but it also exposes enough action/species
information for a simple mechanistic solver to build a strong planning model.

\section{Methods and Adaptations}

The method names in this report are benchmark-native adaptations, not
paper-faithful reproductions of the original deep-RL systems.

\begin{center}
\begin{tabularx}{\linewidth}{p{0.18\linewidth}p{0.25\linewidth}X}
\toprule
method & role in this experiment & adaptation used here \\
\midrule
\code{refplan} & general offline MBRL & posterior ensemble planner over a learned
linear \code{ContinuousDynamicsEnsemble}; plans with particle MPC. \\
\code{bamcts} & general offline MBRL & belief-rooted finite-action tree search over
the same learned linear dynamics family. \\
\code{ogsrl} & general offline MBRL & guarded offline policy with learned linear
dynamics, a kNN support guardian, safety/OOD constraints, and a low-capacity
linear actor. \\
\code{moor_native} & native ecological baseline & single fitted Ricker model,
discrete abundance belief, tabular transition/reward model, and value-iteration
action values. This is the genuinely naive ecological baseline. \\
\code{plus_native} & native ecological solver & posterior over four mechanistic
forms, each discretized and solved independently. This is stronger than a naive
baseline because it can identify the true form from observations. \\
\bottomrule
\end{tabularx}
\end{center}

The distinction between adapted and native ecological methods matters. Earlier
benchmark methods named \code{plus} and \code{moor} are mechanistic
particle-MPC adaptations. The motivation experiment instead uses
\code{plus_native} and \code{moor_native}, routed through
\code{filter=native_discrete}, so the native methods use their own discretized
beliefs and tabular solvers rather than the shared learned particle filter.

\code{moor_native} is naive in the relevant sense: it commits to a single
Ricker-form model even when the true family is Allee, theta-logistic, or
regime-switching. \code{plus_native} is not naive. It maintains a posterior over
the four mechanistic forms. A one-episode diagnostic on Iberian lynx,
$\sobs=0.2$, safe mode shows that \code{plus_native} rapidly identifies the
true form:

\begin{center}
\begin{tabular}{lrrrr}
\toprule
true family & $P(\mathrm{Ricker})$ & $P(\mathrm{Allee})$ & $P(\theta)$ & $P(\mathrm{regime})$ \\
\midrule
Ricker & 0.999 & 0.000 & 0.000 & 0.001 \\
Allee & 0.000 & 0.901 & 0.001 & 0.098 \\
theta-logistic & 0.000 & 0.021 & 0.917 & 0.062 \\
regime-switching & 0.000 & 0.097 & 0.000 & 0.903 \\
\bottomrule
\end{tabular}
\end{center}

Thus the best-native statistic is best understood as a comparison against the
best ecological solver, including one that is close to a form-oracle. The
cleaner naive-baseline result is the comparison against \code{moor_native}
alone.

\section{Why the General Methods Fail Here}

The general methods fail because they are model-limited. They must learn a
dynamics model from the offline dataset, then plan through that learned model.
The native methods instead build mechanistic transition tables from public
ecological structure: action-specific growth set-points, carrying-capacity
effects, and species-scale carrying capacity.

Changing \code{filter=ricker} or \code{filter=true_family} would not give the
general methods the same information. Those settings change state estimation,
but \code{refplan}, \code{bamcts}, and \code{ogsrl} still fit and plan with
their own learned \code{ContinuousDynamicsEnsemble}. A valid information
ablation must change the planning model source itself.

The search-budget audit supports this diagnosis. Increasing planner/search
budget closes only 21\% of the mean gap, and deeper rollouts make all three
general methods worse on recoverable populations. That is the expected pattern
when model error compounds over rollout depth. The failure is therefore not
simply an optimization crash, a missing training log, or an insufficient search
budget.

\section{What Plots Are Available}

\begin{center}
\begin{tabularx}{\linewidth}{p{0.26\linewidth}p{0.35\linewidth}X}
\toprule
plot family & available from completed artifacts? & interpretation \\
\midrule
training / validation loss
& Partly. Every row has \code{training_history.csv/json}; the main run has
1152 per-row training-history PNGs. \code{refplan}, \code{bamcts}, and
\code{ogsrl} record train/holdout dynamics log-MSE; \code{ogsrl} additionally
records an actor-loss trace over update steps.
& Useful for sanity and model-quality diagnostics. It does not rescue the
original motivation: the learned models have finite holdout error, while the
native baselines use nearly exact public ecological tables. \\
\addlinespace
return / reward over time
& Per-episode total return is available in \code{episodes.csv} and
\code{summary.json}. Per-timestep reward traces are not saved by the completed
main evaluator.
& Aggregate return figures are available and are the right primary evidence.
Per-timestep reward curves would require a targeted trace regeneration, not a
new headline experiment. \\
\addlinespace
actions over time, all methods in one figure
& Not as a saved artifact of the completed main run. The evaluator stores action
entropy and danger-zone action fractions, but not the full per-timestep action
sequence. The repository already has trajectory capture/plotting scripts; they
need native-method routing before they can produce the requested all-method
panel for this experiment.
& Do not claim action-over-time evidence from the audited sweep itself. A paired
illustrative trace can be regenerated from frozen configs and datasets, but it
should be labelled illustrative rather than treated as a headline statistic. \\
\addlinespace
action taken in danger zone
& Yes. \code{episodes.csv} includes
\code{danger_action_0_fraction} through \code{danger_action_10_fraction}.
& This is a meaningful safety diagnostic. It should be reported as a
conditional action-composition check, with recoverable populations separated
from demographic sinks. \\
\bottomrule
\end{tabularx}
\end{center}

\section{Headline Returns}

\begin{figure}[t]
\centering
\includegraphics[width=0.82\linewidth]{docs/benchmark/figures/motivation_native/moor_native_returns_by_family.pdf}
\caption{Recoverable safe-mode cells, pooled over observation-noise levels.
Bars show the best general method versus \code{moor_native}, the single-Ricker
native baseline. The gap is negative for every dynamics family and is larger on
the non-Ricker families where the Ricker assumption is false. This is the clean
naive-baseline comparison, without the \code{plus_native} form-identification
confound.}
\label{fig:motivation-moor-family}
\end{figure}

\begin{center}
\begin{tabular}{lrrr}
\toprule
family & best general & \code{moor_native} & gap \\
\midrule
Ricker & 7.749 & 8.386 & -0.636 \\
Allee & 7.762 & 8.544 & -0.782 \\
theta-logistic & 7.828 & 8.777 & -0.950 \\
regime-switching & 7.462 & 8.370 & -0.908 \\
all & 7.700 & 8.519 & -0.819 \\
\bottomrule
\end{tabular}
\end{center}

Figure~\ref{fig:motivation-moor-family} is the cleanest figure for the
motivation. It does not support the original hypothesis. A misspecified
single-Ricker ecological solver beats the general methods by more on the
families where its own assumption is false than on the family where it is
correct. The best general method beats \code{moor_native} in only 6 of 112
recoverable safe-mode cells. Model-form misspecification costs the native
solver less than learned-model error costs the general methods.

\begin{figure}[t]
\centering
\includegraphics[width=0.82\linewidth]{docs/benchmark/figures/motivation_native/headline_returns_by_family.pdf}
\caption{Recoverable safe-mode cells, pooled over observation-noise levels.
Bars show the best general method versus the best native ecological solver,
where the native side may choose between \code{moor_native} and
\code{plus_native}. This is the strongest ecological-solver comparison, but it
should not be described as purely naive because \code{plus_native} identifies
the mechanistic form quickly.}
\label{fig:motivation-family}
\end{figure}

\begin{center}
\begin{tabular}{lrrr}
\toprule
family & best general & best native solver & gap \\
\midrule
Ricker & 7.749 & 8.420 & -0.670 \\
Allee & 7.762 & 8.895 & -1.133 \\
theta-logistic & 7.828 & 8.935 & -1.107 \\
regime-switching & 7.462 & 8.698 & -1.235 \\
\bottomrule
\end{tabular}
\end{center}

Figure~\ref{fig:motivation-family} shows that adding the model-identifying
\code{plus_native} solver widens the gap, but does not create it. The negative
result already holds against the genuinely naive \code{moor_native} baseline.

\begin{figure}[t]
\centering
\includegraphics[width=\linewidth]{docs/benchmark/figures/motivation_native/all_methods_returns_by_family.pdf}
\caption{All five methods on recoverable safe-mode cells, averaged by dynamics
family. The figure shows that the native advantage is not caused by one weak
general method: all three general methods sit below \code{moor_native} and
\code{plus_native} on every family.}
\label{fig:motivation-all-methods}
\end{figure}

\begin{figure}[t]
\centering
\includegraphics[width=0.74\linewidth]{docs/benchmark/figures/motivation_native/native_gap_by_sigma.pdf}
\caption{Best-general minus best-native return on recoverable safe cells by
observation-noise level. The gap remains negative at every $\sobs$. Increased
observation noise is therefore not the source of the native-baseline advantage.}
\label{fig:motivation-sigma}
\end{figure}

\section{Training Diagnostics}

\begin{figure}[t]
\centering
\includegraphics[width=0.74\linewidth]{docs/benchmark/figures/motivation_native/training_dynamics_log_mse.pdf}
\caption{Train and holdout dynamics log-MSE for the learned general methods,
averaged over the main sweep. This is a model-quality diagnostic. It should not
be read as a training-convergence plot for \code{refplan} or \code{bamcts},
which record final fit diagnostics rather than a multi-step optimization trace.}
\label{fig:motivation-dynmse}
\end{figure}

\begin{figure}[t]
\centering
\includegraphics[width=0.74\linewidth]{docs/benchmark/figures/motivation_native/ogsrl_surrogate_loss.pdf}
\caption{\code{ogsrl} actor surrogate loss over update steps, averaged over
the main sweep. The trace is stable enough to rule out an obvious optimization
crash, but stable training does not imply competitive ecological control.}
\label{fig:motivation-ogsrl-loss}
\end{figure}

The training plots are meaningful, but they support a diagnostic story rather
than the original motivation. The general methods are not failing because the
training logs are missing or because \code{ogsrl} visibly diverges. The
search-budget audit is more decisive: increasing planner/search budget closes
only 21\% of the mean gap, and deeper rollouts make all three general methods
worse on recoverable populations. That pattern is consistent with compounding
model error, not with a simple shortage of search.

\begin{figure}[t]
\centering
\includegraphics[width=0.72\linewidth]{docs/benchmark/figures/motivation_native/search_budget_gap.pdf}
\caption{Clean search-budget audit on 72 recoverable safe cells. The as-run
best general method trails the best native solver by 1.216 mean return; the
best tuned general configuration still trails by 0.962. Tuning closes only
21\% of the gap and raises the win count from 2/72 to 3/72, so the negative
result is not explained by insufficient search budget.}
\label{fig:motivation-search-budget}
\end{figure}

\section{Danger-Zone Actions}

\begin{figure}[t]
\centering
\includegraphics[width=\linewidth]{docs/benchmark/figures/motivation_native/danger_zone_action_composition.pdf}
\caption{Conditional action composition when the previous true state is in the
danger band $\ssafe < s_t \le 4\ssafe$, split between recoverable populations
and demographic sinks. The figure is computed from the
\code{danger_action_*_fraction} fields in \code{episodes.csv}. It is an
action-selection diagnostic, not a per-timestep trajectory trace.}
\label{fig:motivation-danger-actions}
\end{figure}

The danger-zone action data are available and useful, especially for checking
that policies do not all collapse to the same passive action near the safety
boundary. They do not, however, rescue the original conservation-motivation
claim. Among recoverable populations, collapse entry is essentially zero for
both general and native methods at every noise level, but the danger band is not
empty: 7{,}099 of 11{,}200 recoverable safe-mode episodes (63.4\%) enter
$\ssafe < s_t \le 4\ssafe$. Thus Figure~\ref{fig:motivation-danger-actions}
does measure decisions under real safety pressure. The distinction is that
recoverable populations repeatedly approach the boundary and avoid collapse,
whereas the two demographic sinks are structurally different; they contribute
1{,}600 of 3{,}200 safe-mode danger-band episodes and are therefore shown
separately.

\section{What the Plots Support}

The figures support the following reading:
\begin{itemize}
  \item The native ecological baselines are strong because the environment gives
  them public per-action growth set-points and public carrying-capacity
  baselines. For \code{moor_native}, remaining Ricker-form misspecification is
  second-order for action ranking; for \code{plus_native}, the form posterior
  resolves the structural family quickly.
  \item The general methods are model-limited. They must learn dynamics from
  4000 logged transitions, and increasing rollout depth worsens performance
  rather than improving it.
  \item The original positive motivation---general offline MBRL beats naive
  ecological baselines under structural uncertainty---is not supported by this
  experiment.
\end{itemize}

The figures therefore support a reframed paper result: in this real-ecology
setting, a well-specified discretized ecological solver is a very strong
baseline, and even a misspecified single-Ricker native solver is hard for
general offline MBRL to beat. That is still a meaningful scientific result, but
it is a negative result relative to the planned story.

\section{Potential Next Step}

The most decision-relevant next step is a planning-model source ablation. The
current interpretation is that the native baselines win because the benchmark
gives them public action-specific growth and carrying-capacity structure, while
the general methods must learn an approximate dynamics model from the offline
dataset. The existing evidence is consistent with this diagnosis, especially
because deeper search worsens the general methods, but it is still indirect.

The direct test is to give the general planners the same mechanistic dynamics
source as the native ecological baselines while keeping the rest of their
planning machinery intact. This cannot be tested by merely changing
\code{filter=ricker} or \code{filter=true_family}: those settings change the
state-estimation proposal, while \code{refplan}, \code{bamcts}, and
\code{ogsrl} still fit and plan with their own learned
\code{ContinuousDynamicsEnsemble}. A valid ablation therefore needs an explicit
planning-model switch, for example
\code{model.dynamics_source \in \{learned,ricker,true_family\}}, with
\code{learned} as the default and a bit-identical no-op proof before any new
run is accepted.

This ablation is the experiment that can still change the paper decision. If
the general methods close most of the native gap when they plan with the same
mechanistic model, the informational-asymmetry diagnosis is confirmed and an
environment redesign is a plausible route to recovering the original
motivation. If the gap remains, the diagnosis must shift from information to
solver structure or policy class, and redesigning only the public ecological
tables is unlikely to rescue the planned story. In that case the honest paper
framing is the negative result: in this real-ecology setting, native ecological
solvers are hard to beat.

A paired trajectory diagnostic remains useful, but as an illustrative companion
rather than the primary next experiment. After the model-source ablation, a
small set of pre-registered cells---for example Amur tiger/Ricker, Iberian
lynx/Allee, and Egyptian vulture/theta-logistic---should be regenerated with
per-timestep true state, observation, action, and reward for all five methods.
Those traces can show whether the measured gap comes from earlier intervention,
more persistent capacity-building actions, avoidance of harvest near the safety
boundary, or different responses to model error. They should remain qualitative
mechanism figures; the audited quantitative claims should remain the 1440-row
main sweep, the 1152-row resolution replicate, the 1728-row search-budget
audit, and any future model-source ablation.

\section{If Per-Timestep Figures Are Needed}

Reward-over-time and action-over-time figures for all methods in one panel
cannot be derived from the completed main-run artifacts alone. The evaluator
kept the per-step arrays internally but saved only per-episode summaries. The
repository already contains \path{scripts/capture_trajectories.py} and
\path{scripts/plot_trajectories.py}; the missing work is to route
\code{plus_native}/\code{moor_native} through \code{native_discrete} and reuse
the frozen real-ecology datasets rather than the synthetic default. Those
traces should be labelled illustrative, because the audited statistics remain
the 1440-row main sweep and the 1152-row resolution replicate.

\section*{Verification}

Verified against: \path{real_ecology_runs/motivation_native_20260711/analysis/aggregate_main.json}, \path{real_ecology_runs/motivation_native_20260711/analysis/resolution_check.json}, \path{real_ecology_runs/motivation_native_20260711/outputs/main}, \path{real_ecology_runs/general_audit_20260712/analysis/budget_verdict_clean.json}, \path{docs/benchmark/SERVER_REVIEW_standalone_report_readiness.md}, \path{src/real_ecology_benchmark/evaluator.py}, \path{scripts/plot_motivation_native_results.py} @ e74526e5e70d959d02311df8f095479f937e1d7b

\end{document}
